\documentclass[11pt]{report}
\usepackage[T1]{fontenc}
\usepackage[margin=30mm]{geometry}
\setcounter{secnumdepth}{3}
\setcounter{tocdepth}{3}
% Fixture convention: one standard heading or caption with a label per line.
% Plain titles deliberately contain no TeX markup or escaped characters.
\newcommand{\NetworkDiagram}[3]{%
  \setlength{\unitlength}{1mm}%
  \begin{picture}(100,35)
    \put(5,28){\circle*{2}}\put(5,28){\vector(2,-1){38}}
    \put(5,5){\circle*{2}}\put(5,5){\vector(2,1){38}}
    \put(48,17){\circle*{3}}\put(48,17){\vector(1,0){42}}
    \put(8,30){\makebox(0,0)[l]{#1}}
    \put(8,2){\makebox(0,0)[l]{#2}}
    \put(65,22){\makebox(0,0)[l]{#3}}
    \put(92,17){\circle{4}}
  \end{picture}%
}
\title{Listening to Imaginary Rivers\\Sensor Placement in Simulated Branching Networks}
\author{Mira Fenwick\\Synthetic Research Design Laboratory}
\date{October 2026\\Original fictional thesis fixture}
\begin{document}
\maketitle
\begin{abstract}
This fictional research thesis develops a reproducible thought experiment for placing environmental sensors in simulated river networks. A small directed network represents tributaries, confluences, and downstream travel. Candidate sensors are compared using coverage, delay, and resilience to missing observations. Every network, measurement, and numerical result is invented for this document; none describes an actual river, field campaign, institution, or published study. The purpose is to connect a coherent research narrative to inspectable source files and a local source-navigation interface. The proposed comparison suggests, within its illustrative assumptions only, that protecting distinct upstream branches can be more useful than concentrating sensors near a single outlet. The thesis separates such conditional observations from claims that would require real calibration and field evidence.
\end{abstract}
\tableofcontents
\listoffigures
\listoftables
\chapter{Introduction}\label{demo-framing}
A sensor can only report the water that reaches it. In a branching river, this simple fact creates a placement problem: a downstream instrument may observe many sources but conceal where a disturbance began. This chapter defines the imagined problem without borrowing a field site or a real monitoring record.
\section{Research questions and objectives}\label{demo-framing-question}
We ask how a small instrument budget should be distributed across branches when disturbances travel downstream and observations can fail. The target is a transparent decision procedure, not a prediction for an existing watershed.
\subsection{Scope of the tracer model}\label{demo-framing-scope}
The experiment concerns a dissolved tracer with an arbitrary unit of concentration. No toxicity threshold is implied. All candidate sites are assumed accessible and equally costly unless a scenario explicitly changes that assumption.
\subsection{Event detection criteria}\label{demo-framing-evidence}
An event is detected when a simulated concentration exceeds a fixed threshold at an available sensor. We record event coverage separately from arrival delay so that late detection cannot masquerade as prompt warning.
\section{Representing downstream transport through branching networks with shared observation sites}\label{demo-framing-network}
Edges point downstream. A node may receive two upstream flows, but no edge returns water to an ancestor. This excludes tides and recirculating reservoirs by construction.
\begin{figure}[htbp]
\centering
\NetworkDiagram{West branch}{East branch}{Outlet}
\caption{An invented two-branch river, with a confluence and one downstream observation site.}\label{demo-fig-network}
\end{figure}
\subsection{Synthetic measurement units}\label{demo-framing-units}
Distance is measured in synthetic reaches and time in ticks. A reach has no fixed conversion to kilometres; a tick has no fixed conversion to hours. These choices prevent unwarranted geographic interpretation.
\begin{table}[htbp]
\centering
\begin{tabular}{lll}
\hline
Symbol & Meaning & Illustrative range \\
\hline
$n$ & Candidate sites & 7--31 \\
$b$ & Sensor budget & 1--5 \\
$t$ & Arrival tick & 0--24 \\
\hline
\end{tabular}
\caption{Symbols used throughout the fictional experiment.}\label{demo-tab-symbols}
\end{table}
\section{Scope and limitations}\label{demo-framing-boundaries}
An efficient arrangement on a graph is not necessarily a feasible deployment. Power, access permissions, maintenance, and habitat disturbance remain outside the simulated objective.
\subsection{Consistency checks}\label{demo-framing-checks}
The first check is conceptual: every numerical statement must identify the simulated world that produced it. Where the thesis uses a hand-worked example rather than a computed trial, that distinction is made explicit.

\chapter{Imaginary catchments: Constructing branching networks with controlled transport and release conditions}\label{demo-worlds}
The simulated worlds vary one structural property at a time. A balanced tree contrasts with a long main stem receiving short tributaries. Neither is intended as a map of a particular place.
\section{Generating contrasting catchments with comparable source counts and different travel paths}\label{demo-worlds-topology}
Each family begins with an outlet and adds upstream branches under a fixed depth limit. Node identifiers are local labels rather than geographic coordinates. Relabelling nodes must leave the physical interpretation unchanged.
\subsection{Balanced tributary networks}\label{demo-worlds-balanced}
A balanced network gives comparable path lengths to most headwaters. It is useful for isolating the value of branch diversity because no tributary dominates travel time by design.
\begin{figure}[htbp]
\centering
\NetworkDiagram{Equal length}{Equal length}{Shared reach}
\caption{Balanced tributaries isolate the effect of sharing a downstream reach.}\label{demo-fig-balanced}
\end{figure}
\subsection{Asymmetric travel paths along an extended main stem}\label{demo-worlds-stem}
The contrasting network places most travel distance along one stem. Short side branches arrive at different positions, allowing a sensor near the outlet to have broad coverage but uneven warning times.
\section{Transport and mixing assumptions}\label{demo-worlds-transport}
A pulse moves one reach per tick in the simplest world. At a confluence, incoming concentrations are combined using equal weights. The model intentionally omits dispersion and chemical reactions.
\subsection{Separating source releases across non-overlapping trials}\label{demo-worlds-pulses}
One pulse is released at each eligible source, with trials separated sufficiently to avoid overlap. This schedule tests source coverage rather than the separation of simultaneous disturbances.
\begin{table}[htbp]
\centering
\begin{tabular}{lrrr}
\hline
World & Sites & Maximum depth & Pulses \\
\hline
Balanced miniature & 7 & 3 & 4 \\
Stem miniature & 9 & 6 & 4 \\
Wide miniature & 13 & 4 & 8 \\
\hline
\end{tabular}
\caption{Hand-specified miniature worlds; values are illustrative design choices, not sampled river statistics.}\label{demo-tab-worlds}
\end{table}
\subsubsection{Final ticks}\label{demo-worlds-boundary}
A pulse arriving exactly at the final observation tick is counted as observed. This inclusive convention matters for long paths and is held constant across all placement methods.
\section{Consistency checks}\label{demo-worlds-checks}
A valid world has exactly one outlet and a downstream path from every source to that outlet. The miniature examples can be inspected by tracing arrows; larger worlds would require the same conditions to be checked mechanically.

\chapter{Sensor placement under constrained budgets: Comparing greedy coverage and outage-aware selection in branching networks}\label{demo-placement}
Placement is treated as a sequence of explicit choices. We compare a downstream baseline with a greedy coverage rule, then examine what changes when unavailable sensors are part of the design rather than a surprise at evaluation time.
\section{Defining sensor coverage without conflating detection delay and source attribution}\label{demo-placement-objective}
For a set of sensors, coverage is the fraction of scheduled source pulses observed at least once before the final tick. Multiple detections of one pulse do not increase that fraction.
\subsection{Balancing warning delay and branch coverage when downstream readings cannot distinguish source events}\label{demo-placement-coverage}
A broad footprint can be useful even when attribution is poor. We therefore keep coverage and warning delay as separate outputs rather than hiding their tension in an unexplained weighted score. Simultaneous pulses are discussed as a limitation, not evaluated as if the single-pulse experiment had resolved them.
\subsubsection{Resolving equal coverage gains through warning delay and a stable candidate ordering}\label{demo-placement-ties}
When two candidate sites add the same coverage, the rule chooses the one with smaller mean arrival delay. A remaining tie is resolved by the fixed local node order. This last rule ensures a repeatable example, not a scientific preference for a label.
\begin{table}[htbp]
\centering
\begin{tabular}{lrrl}
\hline
Candidate & Added pulses & Mean delay & Decision \\
\hline
Upper fork & 2 & 2 & Select \\
Lower fork & 2 & 4 & Defer \\
Outlet & 1 & 5 & Defer \\
\hline
\end{tabular}
\caption{An illustrative greedy step: delay resolves a coverage tie.}\label{demo-tab-ties}
\end{table}
\section{Baselines}\label{demo-placement-baselines}
The outlet-only baseline spends one sensor at the last node. A spaced-stem baseline fills the budget along the main stem at approximately equal graph distances. Both are intentionally simple enough to inspect without specialised optimisation software.
\subsection{Incremental sensor selection}\label{demo-placement-greedy}
The greedy rule evaluates each remaining candidate against the currently selected set, adds the best candidate, and repeats until the budget is exhausted. A selected site cannot be selected again. The procedure makes no claim of global optimality under the combined tie-breaking rule.
\begin{figure}[htbp]
\centering
\fbox{\begin{minipage}{0.8\linewidth}
Start with no sensors.\\[4pt]
Compare the additional covered pulses at each unused site.\\[4pt]
Choose by coverage, then delay, then local order.\\[4pt]
Stop when the fixed budget is spent.
\end{minipage}}
\caption{The explicit selection sequence used in the worked examples.}\label{demo-fig-selection}
\end{figure}
\subsection{Accounting for sensor outages in the placement objective}\label{demo-placement-outages}
A resilience variant scores each candidate across a small list of availability masks. Each mask removes a sensor reading rather than deleting a river node, so water transport remains unchanged.
\subsubsection{Outage masks}\label{demo-placement-masks}
The masks include all sensors available and each single-sensor outage. Equal mask weights are an illustrative assumption. They are not estimated failure probabilities and should not be read as reliability measurements.
\section{Computational cost of comparing candidate sites across repeated placement steps}\label{demo-placement-cost}
For each placement step, the rule compares unused sites against a fixed event list. Precomputed arrival relationships avoid repeating the graph walk for every candidate. This observation concerns the proposed algorithm; no benchmark timings are claimed here.

\chapter{Illustrative comparisons: Separating source coverage from observation delay when sensor readings are unavailable}\label{demo-evaluation}
This chapter uses hand-authored numbers to demonstrate reporting structure. The values are internally described examples, not outputs from an implemented simulator, experiments on real water, or estimates of field performance.
\section{Experimental comparison protocol}\label{demo-evaluation-protocol}
Each comparison holds the world, pulse schedule, and budget fixed. Only the placement rule changes. A separate comparison changes the world while retaining the rule; mixing these two operations would obscure the reason for a difference.
\subsection{Coverage and delay metrics}\label{demo-evaluation-metrics}
Coverage counts observed pulses out of four scheduled pulses in the miniature example. Mean delay is reported only for observed pulses. A method missing most pulses must not be praised solely for a low conditional delay.
\subsubsection{Denominators}\label{demo-evaluation-denominators}
The count of observed pulses accompanies every coverage fraction. Unobserved pulses have no finite observed delay, so they remain visible in the count rather than being silently removed from the comparison.
\begin{table}[htbp]
\centering
\begin{tabular}{lrrr}
\hline
Rule, budget two & Observed pulses & Coverage & Mean delay \\
\hline
Spaced stem & 3 of 4 & 0.75 & 4.0 \\
Greedy branches & 4 of 4 & 1.00 & 3.0 \\
Resilience variant & 4 of 4 & 1.00 & 3.5 \\
\hline
\end{tabular}
\caption{Invented reporting example for four pulses: coverage and conditional delay are shown together.}\label{demo-tab-results}
\end{table}
\section{Evaluating downstream backup observations when upstream sensor readings become unavailable}\label{demo-evaluation-outages}
Removing the sole reading from one branch can reveal whether another sensor provides a downstream opportunity. The imagined outage comparison holds the pulse schedule fixed and changes only availability.
\begin{figure}[htbp]
\centering
\NetworkDiagram{Unavailable site}{Available site}{Backup reading}
\caption{Illustrative outage scenario in which a missing upstream reading leaves a downstream backup available, preserving event coverage while delaying the first observation and reducing the ability to attribute a pulse to its original tributary.}\label{demo-fig-outages}
\end{figure}
\subsection{Recovery}\label{demo-evaluation-recovery}
Recovery here means retaining at least one observation of a pulse. It does not mean repairing an instrument or reconstructing its missing concentration series. This narrow definition prevents a coverage result from implying a maintenance capability.
\section{Consistency checks}\label{demo-evaluation-checks}
The example fractions agree with their counts, and the reported delays use only observed pulses. These are necessary reporting checks but do not validate the transport model against reality.
\subsection{Sensitivity to observation-window boundaries}\label{demo-evaluation-sensitivity}
Shortening the observation window can remove late detections without changing any sensor location. A future implementation should report this sensitivity separately from outage sensitivity because the two mechanisms have different interpretations.
\subsubsection{Candidate-order sensitivity}\label{demo-evaluation-order}
If candidate labels are permuted, only exact ties should change the selected set. A change outside a tie would signal dependence on bookkeeping rather than on the stated objective.

\chapter{Interpreting sensor placement under simplified river transport assumptions}\label{demo-interpretation}
The miniature worlds support reasoning about the objective, not general conclusions about river monitoring. Their main value is exposing trade-offs that a single aggregate score might conceal.
\section{Consequences of distributing sensors across separate tributaries instead of a shared downstream reach}\label{demo-interpretation-diversity}
Placing instruments on separate tributaries can reveal different source events earlier than a shared downstream instrument. The benefit depends on the observation window and on whether attribution matters to the intended decision.
\subsection{The cost of delayed observations}\label{demo-interpretation-delay}
A downstream backup trades warning time for redundancy. In the fictional model, both quantities are measurable. In a real deployment, the importance of a delay would depend on a response process that this document does not model.
\begin{figure}[htbp]
\centering
\fbox{\begin{tabular}{ccc}
Upstream observation & $\longrightarrow$ & Earlier warning \\
Shared downstream observation & $\longrightarrow$ & Broader backup
\end{tabular}}
\caption{Two different benefits, not a universal ranking of sensor locations.}\label{demo-fig-tradeoff}
\end{figure}
\subsubsection{Warnings after the response deadline}\label{demo-interpretation-window}
A warning arriving after a hypothetical response deadline may still contribute to retrospective diagnosis. The intended use therefore determines whether the same observation is considered timely, useful but late, or irrelevant.
\section{Limits of the transport model}\label{demo-interpretation-limits}
Constant travel times omit seasonal variation. Equal mixing weights omit unequal discharge. Independent availability masks omit shared power or communications failures. Each omission makes the example easier to inspect and less representative of a deployed network.
\subsection{Mixing assumptions}\label{demo-interpretation-mixing}
A real confluence can produce spatially uneven concentration. The single-value node representation assumes complete mixing immediately. Relaxing this assumption would require more than changing a caption or adding another candidate site.
\begin{table}[htbp]
\centering
\begin{tabular}{ll}
\hline
Simplification & Evidence needed before application \\
\hline
Fixed travel time & Repeated flow observations \\
Equal site costs & Access and maintenance assessment \\
Independent outages & Infrastructure dependency map \\
\hline
\end{tabular}
\caption{Missing evidence that separates an illustrative graph exercise from deployment advice.}\label{demo-tab-limits}
\end{table}
\section{Reporting assumptions and results}\label{demo-interpretation-reporting}
A useful report states which events were included, which readings were unavailable, and which assumptions determined arrival. The original fictional numbers here are labelled at their point of use to avoid implying an unperformed computational study.
\subsection{Consistency checks}\label{demo-interpretation-checks}
A reader should be able to distinguish a design choice, a mathematical consequence, and an empirical claim. This document supplies examples of the first two; it deliberately supplies no empirical river findings.

\chapter{Conclusions}\label{demo-outlook}
The thesis frames sensor placement as a transparent comparison among limited choices. Its contribution is an original worked narrative and inspectable synthetic document, not a validated monitoring system.
\section{Synthesis}\label{demo-outlook-summary}
Coverage, warning delay, and resilience answer different questions. Keeping them separate makes the consequences of a placement rule easier to discuss and avoids treating a downstream observation as equivalent to an early upstream warning.
\subsection{Conclusions within the model}\label{demo-outlook-answer}
Under the imagined branching topology, spreading sensors can protect distinct source paths. A shared downstream sensor can provide backup. Neither strategy dominates without specifying the event schedule, observation window, and outage assumptions.
\section{From illustrative networks to a calibrated environmental monitoring study}\label{demo-outlook-next}
A genuine study would first implement the specified event accounting and verify it against hand-worked networks. Only then would expanded topology families and uncertain travel times provide meaningful additional comparisons.
\subsection{Independent transport calibration}\label{demo-outlook-calibration}
Connecting the model to a real catchment would require independent measurements of transport and mixing. Choosing a parameter because it improves a placement result would confound calibration with the desired conclusion.
\subsubsection{Separating calibration from placement comparisons}\label{demo-outlook-separation}
The data used to choose model parameters should be distinguished from the scenarios used to compare placement rules. This separation makes it easier to identify which conclusion depends on which assumption.
\begin{figure}[htbp]
\centering
\fbox{\begin{minipage}{0.8\linewidth}
Inspect miniature cases $\longrightarrow$ implement event accounting\\[4pt]
Challenge assumptions $\longrightarrow$ gather independent evidence\\[4pt]
Evaluate feasibility $\longrightarrow$ consider a deployment decision
\end{minipage}}
\caption{A conditional path beyond this fictional document; none of these later activities is claimed as completed research.}\label{demo-fig-outlook}
\end{figure}
\section{Closing remarks}\label{demo-outlook-closing}
Listening to an imaginary river is useful when it clarifies what a real experiment would need to establish. The boundary between an inspectable example and a supported claim is part of the research design, not an afterthought.

\appendix
\chapter{Event ledger}\label{demo-appendix}
This appendix supplies a hand-worked ledger for an imaginary three-node chain. It illustrates the bookkeeping convention without claiming to reproduce the separate illustrative comparison table.
\section{A three-node transport example}\label{demo-appendix-chain}
Node One drains to Node Two, which drains to Node Three. A pulse starts at Node One at tick zero and moves one edge each tick. There is no attenuation in this miniature case.
\subsection{Recording arrivals at successive observation sites}\label{demo-appendix-arrivals}
A sensor at Node Two records the pulse at tick one. A sensor at Node Three records it at tick two. With both readings available the first-observation delay is one tick, not the average of the two arrival ticks.
\begin{table}[htbp]
\centering
\begin{tabular}{rlll}
\hline
Tick & Pulse location & Node Two sensor & Node Three sensor \\
\hline
0 & Node One & No arrival & No arrival \\
1 & Node Two & Arrival & No arrival \\
2 & Node Three & No new arrival & Arrival \\
\hline
\end{tabular}
\caption{A literal three-tick ledger for one invented pulse.}\label{demo-tab-ledger}
\end{table}
\subsubsection{Missing first observations}\label{demo-appendix-missing}
If the Node Two sensor is unavailable, the first observed arrival moves to tick two. Coverage remains one out of one when the final tick is two, but becomes zero out of one when the final tick is one.
\section{Document provenance}\label{demo-appendix-provenance}
All prose, diagrams, identifiers, and tables in this fixture are newly authored synthetic material. The source and compiled reading copy demonstrate project-wide navigation with the LaTeX TOC extension. Its outline is derived from source commands, not from the compiled PDF, and does not require a TeX build.

\end{document}
